\scanpage{19}
\margin{\unsure{1}\quad\unsure{?}}
\deleted{\unsure{Revenons|Remarque} : \unsure{à la} question de dualité. Cela \unsure{veut} \unsure{dire}
la validité du (iii) du lemme \unsure{3|2} signifie
que \unsure{la dimension} $\dim A\leq n$. \unsure{et que} \illegible{2}
\unsure{trouve que|donne}}

\uline{Théorème 4}

\uline{Définition} Soit $A$ anneau \deleted{local} noethérien de dim $n$.
Un module $\Omega$ sur $A$ est dit un \textquotedblleft module fondamental\textquotedblright{}
\deleted{\unsure{coefficient de dualité}} si il satisfait les deux conditions

% slanted
\margin{Introduire $\Omega$ par une proposition, disant que (i) et (ii) équivalent à certaines propriétés du foncteur $M\to\operatorname{Hom}(M,\Omega)$.}

(i) $\Omega$ est de type fini

(ii) $\deleted{\unsure{H^n(\Omega)}} D(\Omega)\simeq H^n(A)$

Un tel $\Omega$ est unique à isom. près \deleted{\unsure{au|si}} \inserted{\uline{si}}
$A$ est local (ou semi-local), \unsure{puis que|puisque|mais que} $\deleted{\unsure{D(A)}=}\inserted{\widehat\Omega=D(H^n(A))}$
est alors \unsure{connu|unique}. On le note $\Omega_A$. D'autre part si $A$ est \inserted{local} complet, un tel $\Omega_A$ existe
\unsure{tjrs} (on prend $D(H^n(A))$). Il existe aussi
si $A$ régulier \inserted{et équidimensionnel} : prendre $\Omega=A$ !

\deleted{D'ailleurs, D'après le \unsure{th'|th.} 3,
l'existence de $\Omega$ implique que \unsure{est} \unsure{lgs de un seul}
\illegible{1} \unsure{annulateur|accouplement} de \unsure{ce} $H^n(A)=$}

\margin{Si $A$ est artinien, alors \uline{module fond} $=$ \uline{module dualisant}}

\deleted{\uline{Proposition} Soit $A$ un anneau local}

\uline{Théorème 4} Soit $A$ \inserted{\unsure{locl|local}} noethérien\inserted{de dim $n$}, muni d'un
module \deleted{dualisant} \inserted{\uline{fondamental}} $\Omega$. Le \unsure{.} \unsure{un} \deleted{\unsure{isom.}}\inserted{hom}
\[
H^n(\Omega)\simeq H^n(A)\otimes\Omega\xrightarrow{\unsure{g|\eta}} \mathrm I
\]
[déduit d'un isom $H^n(A)\xrightarrow{\sim}D(\Omega)$], définit
des \unsure{accouplements|accouplement} \unsure{donnant} $\unsure{\alpha_i}$, qui est un
isom en dim $n$. Pour que ce soit un
% archive sheet number: 18

